% sections/bilan.tex
% Phase 10: Bilan -- Conclusion et Perspectives (RAPP-07, RAPP-08)

\section{Bilan}
\label{sec:bilan}

\subsection{Conclusion}
\label{subsec:conclusion}

Les 401\,709 transferts WETH du 5 ao\^ut 2024 sur Polygon ont \'et\'e
analys\'es sous six angles compl\'ementaires. Le Tableau~\ref{tab:synthese_rq}
r\'esume les r\'esultats.

\begin{table}[H]
\centering
\caption{Synth\`ese des r\'esultats par question de recherche}
\label{tab:synthese_rq}
\begin{tabular}{llll}
\toprule
\textbf{RQ} & \textbf{M\'etrique cl\'e} & \textbf{R\'esultat}
  & \textbf{Interpr\'etation} \\
\midrule
RQ1 & Densit\'e & $2{,}41 \times 10^{-4}$
  & R\'eseau creux, structure hub \\
RQ2 & Centralit\'e & Top-10 multi-mesures
  & Routeurs DEX dominants \\
RQ3 & Modularit\'e $Q$ & 0{,}6445
  & Compartimentalisation DeFi \\
RQ4 & Pics d'activit\'e & 01h / 13h UTC
  & Liquidations vs panique \\
RQ5 & Sigma $\sigma$ & 115{,}90
  & Petit-monde confirm\'e \\
RQ6 & Gini $G$ & 0{,}9780
  & Concentration extr\^eme \\
\bottomrule
\end{tabular}
\end{table}

Le r\'eseau est creux ($d = 2{,}41 \times 10^{-4}$) mais presque enti\`erement
connect\'e (WCC \`a 97{,}3\,\%). Il s'organise autour de quelques hubs (routeurs
DEX, ponts, coffres-forts de protocoles) qui canalisent l'essentiel des flux
(RQ1, RQ2). Ces hubs structurent 174 communaut\'es ($Q = 0{,}6445$), qui
correspondent aux diff\'erents protocoles DeFi (RQ3). L'analyse temporelle
(RQ4) fait appara\^itre deux pics distincts~: liquidations automatiques \`a
gros volume autour de 01h~UTC, puis panique retail \`a 13h~UTC \`a
l'ouverture de Wall Street. Les m\'etriques petit-monde sortent tr\`es
au-dessus du r\'ef\'erent al\'eatoire ($\sigma = 115{,}90$, RQ5), et la
concentration de volume est extr\^eme ($G = 0{,}9780$, RQ6)~: 1\,\% des
ar\^etes transportent 81{,}3\,\% du volume.

Mis bout \`a bout, ces six r\'esultats dessinent un m\^eme r\'eseau~: rapide et
efficace en temps normal parce que peu de hubs traitent beaucoup de flux,
fragile en temps de crise pour la m\^eme raison. Quelques hubs concentrent
connexions et volumes, les chemins sont courts donc les chocs voyagent vite,
et la liquidit\'e d\'epend d'un petit nombre de pools. Le 5 ao\^ut 2024, ces
propri\'et\'es ont jou\'e ensemble.

\subsection{Perspectives}
\label{subsec:perspectives}

Plusieurs pistes prolongeraient ce travail.

\paragraph{Fitting power-law formel.}
L'analyse topologique (RQ1) a sugg\'er\'e une distribution des degr\'es \`a
queue lourde sans toutefois fournir de test statistique formel. Une analyse
rigoureuse selon la m\'ethodologie de Clauset et
al.~\cite{clauset2009powerlaw} (estimation de l'exposant $\alpha$ et du
seuil $x_{\min}$ par maximum de vraisemblance, suivie d'un test de
Kolmogorov-Smirnov comparant les mod\`eles power-law, lognormal et
exponentiel) permettrait de statuer sur la nature sans \'echelle du
r\'eseau.

\paragraph{Analyse multi-tokens.}
La pr\'esente \'etude se limite aux transferts WETH. \'Etendre l'analyse
aux transferts ETH, USDC et MATIC sur la m\^eme journ\'ee r\'ev\'elerait si
les propri\'et\'es structurelles observ\'ees (topologie hub-and-spoke,
petit-monde, concentration extr\^eme) sont sp\'ecifiques au WETH ou
caract\'eristiques de l'ensemble de l'\'ecosyst\`eme DeFi sur Polygon.

\paragraph{Analyse multi-jours.}
Le clich\'e unique du 5 ao\^ut 2024 capture le crash mais pas la
dynamique de r\'ecup\'eration. Analyser les jours pr\'ec\'edant, pendant et
suivant le Lundi Noir permettrait d'observer l'\'evolution de la topologie
du r\'eseau \`a travers les phases de crise~: stabilit\'e pr\'e-choc,
propagation du choc et restructuration post-choc.

\paragraph{Visualisation interactive.}
On a d\'evelopp\'e en parall\`ele du rapport un prototype web pour explorer
les m\^emes donn\'ees \`a la souris~: \url{https://crypto-graph-explorer.vercel.app}~\cite{cryptographexplorer_web}.
La visualisation est faite en D3.js, l'interface en React. Les six RQ et les
cinq analyses compl\'ementaires de la section~\ref{sec:complementaires} y
sont accessibles depuis un panneau lat\'eral, avec zoom, drag et filtrage
par communaut\'e sur les $12\,880$ sommets. Le prototype tourne dans le
navigateur, donc sur des m\'etriques pr\'e-calcul\'ees et quelques
approximations~: \textit{betweenness} \'echantillonn\'ee, Gini calcul\'e sur
les degr\'es au lieu des volumes WETH, et la dimension temporelle est
absente faute de \textit{timestamps} dans le CSV embarqu\'e. Les chiffres
exacts du rapport ($\sigma$, $Q$, $G$) restent ceux obtenus localement avec
Python. Restent \`a embarquer dans le prototype les \textit{timestamps}
(RQ4) et les poids WETH par ar\^ete (RQ6), et \`a y brancher un
\textit{fitting} formel de la loi de puissance pour la distribution des
degr\'es.
